\documentclass{article}

\usepackage[utf8]{inputenc}
\usepackage[T1]{fontenc}
\usepackage{amsmath,amssymb}
\usepackage{graphicx}
\usepackage{booktabs}
\usepackage{hyperref}
\usepackage{subcaption}
\usepackage[margin=1in]{geometry}

\title{Supplementary Materials:\\Latent Neural SDEs with Conditional Score Matching\\for Probabilistic Solar Irradiance Nowcasting}

\author{Keshav Krishnan}

\begin{document}

\maketitle

\appendix

% ============================================================
\section{Full Hyperparameter Tables}
\label{app:hyperparams}
% ============================================================

\begin{table}[h]
\centering
\caption{CS-VAE hyperparameters.}
\begin{tabular}{lcc}
\toprule
Parameter & Default & Ablated Values \\
\midrule
Latent dimension $d_z$ & 64 & \{16, 32, 128, 256\} \\
$\beta$ (KL weight) & 0.1 & \{0.01, 1.0, 4.0\} \\
Encoder channels & [32, 64, 128, 256, 512] & --- \\
Learning rate & $1 \times 10^{-4}$ & --- \\
Batch size & 64 & --- \\
Epochs & 100 & --- \\
Optimizer & Adam & --- \\
Image size & $256 \times 256$ & --- \\
\bottomrule
\end{tabular}
\end{table}

\begin{table}[h]
\centering
\caption{Neural SDE hyperparameters.}
\begin{tabular}{lcc}
\toprule
Parameter & Default & Ablated Values \\
\midrule
Drift hidden dim & 256 & --- \\
Drift architecture & 2 Linear + 2 ResBlock & --- \\
Diffusion hidden dim & 64 & --- \\
Covariate dimension $d_c$ & 5 & \{0\} (ablation A7) \\
$\lambda_\sigma$ (diffusion loss weight) & 1.0 & \{0.1, 0.5, 5.0\} \\
CTI window size $W$ & 10 frames (100s) & --- \\
SDE step size $\Delta t$ & 10 seconds & --- \\
Learning rate & $1 \times 10^{-4}$ & --- \\
Batch size & 128 & --- \\
Epochs & 200 & --- \\
Gradient clipping & 1.0 & --- \\
\bottomrule
\end{tabular}
\end{table}

\begin{table}[h]
\centering
\caption{Score-Matching Decoder (CSMID) hyperparameters.}
\begin{tabular}{lcc}
\toprule
Parameter & Default & Ablated Values \\
\midrule
Score network hidden dim & 256 & --- \\
Number of residual blocks & 2 & --- \\
Diffusion steps $S$ & 100 & \{50, 200\} \\
$\beta$ schedule & Linear & --- \\
$\beta_{\text{start}}$ & $1 \times 10^{-4}$ & --- \\
$\beta_{\text{end}}$ & 0.02 & --- \\
Learning rate & $1 \times 10^{-4}$ & --- \\
Batch size & 256 & --- \\
Epochs & 100 & --- \\
\bottomrule
\end{tabular}
\end{table}

\begin{table}[h]
\centering
\caption{Inference hyperparameters.}
\begin{tabular}{lcc}
\toprule
Parameter & Default & Ablated Values \\
\midrule
Monte Carlo samples $N$ & 100 & \{10, 25, 50, 200, 500, 1000\} \\
SDE solver & Euler-Maruyama & --- \\
Score decoder reverse steps & 100 & --- \\
Prediction interval level & 90\% & --- \\
\bottomrule
\end{tabular}
\end{table}

% ============================================================
\section{Dataset Details}
\label{app:data}
% ============================================================

\subsection{CloudCV Dataset}

The CloudCV dataset~\citep{perrsauer2024cloudcv} was collected at the NREL Solar Radiation Research Laboratory (SRRL) in Golden, Colorado (39.742\textdegree N, 105.18\textdegree W, 1829m altitude). Sky images were captured using an ELP 180-degree fisheye lens USB camera at 1920$\times$1080 resolution every 10 seconds during daylight hours. Co-located irradiance measurements were taken with a LI-COR LI-200 pyranometer digitized via an Adafruit ADS1115 ADC.

Ground-truth GHI was obtained from the SRRL Baseline Measurement System (BMS) calibrated LI-200 pyranometer, which provides 1-minute resolution measurements. BMS GHI was interpolated to 10-second resolution via linear interpolation to match the CloudCV image timestamps.

\subsection{Data Processing}

\begin{enumerate}
    \item \textbf{Image preprocessing}: Fisheye images were center-cropped to a square region and resized to $256 \times 256$ pixels, with pixel values normalized to $[0, 1]$.
    \item \textbf{Clear-sky computation}: Ineichen-Perez clear-sky GHI was computed for each timestamp using \texttt{pvlib}~\citep{holmgren2018pvlib} with the SRRL location parameters.
    \item \textbf{Clear-sky index}: $k_t = \text{GHI}_{\text{measured}} / \text{GHI}_{\text{clear}}$, clipped to $[0, 1.5]$.
    \item \textbf{Nighttime filtering}: All data with solar zenith angle $> 85$\textdegree\ removed.
    \item \textbf{Ramp detection}: $|\Delta\text{GHI}| > 50\;\text{W/m}^2$ within 60 seconds.
    \item \textbf{Chronological split}: No temporal shuffling; train/val/test are contiguous date ranges.
\end{enumerate}

% ============================================================
\section{Extended Ablation Results}
\label{app:ablation}
% ============================================================

Table~\ref{tab:ablation_full} presents the full ablation results across all forecast horizons. The CTI gating mechanism (A2) shows increasing importance at longer horizons, where cloud dynamics dominate persistence effects. The deterministic ODE ablation (A5) produces the most dramatic degradation in PICP, confirming that stochastic modeling is essential for calibrated uncertainty quantification.

\begin{table}[h]
\centering
\caption{Extended ablation: CRPS across all forecast horizons.}
\label{tab:ablation_full}
\begin{tabular}{lccccccc}
\toprule
Variant & 1 min & 2 min & 5 min & 10 min & 15 min & 20 min & 30 min \\
\midrule
A1: Full & -- & -- & -- & -- & -- & -- & -- \\
A2: $-$CTI & -- & -- & -- & -- & -- & -- & -- \\
A3: $-$VAE & -- & -- & -- & -- & -- & -- & -- \\
A4: $-$Score & -- & -- & -- & -- & -- & -- & -- \\
A5: $-$SDE & -- & -- & -- & -- & -- & -- & -- \\
A6: Adjoint & -- & -- & -- & -- & -- & -- & -- \\
A7: $-$Meteo & -- & -- & -- & -- & -- & -- & -- \\
\bottomrule
\end{tabular}
\end{table}

% ============================================================
\section{Sampling Efficiency}
\label{app:sampling}
% ============================================================

Figure~\ref{fig:sampling} shows CRPS convergence as a function of Monte Carlo sample count $N$. CRPS stabilizes within 1\% of the $N=1000$ reference at approximately $N=50$, confirming that the model produces computationally practical forecasts for real-time deployment. At $N=50$ with a 30-minute horizon, total inference latency remains well under 1 second on GPU hardware.

% ============================================================
\section{Computational Cost Analysis}
\label{app:compute}
% ============================================================

\begin{table}[h]
\centering
\caption{Computational cost breakdown (single GPU, RTX 3090 equivalent).}
\begin{tabular}{lccc}
\toprule
Component & Parameters & Training Time & Epochs \\
\midrule
CS-VAE & 7.2M & $\sim$4 hours & 100 \\
Neural SDE & 0.3M & $\sim$2 hours & 200 \\
Score Decoder & 0.05M & $\sim$20 min & 100 \\
Fine-tuning & 7.55M & $\sim$3 hours & 50 \\
\midrule
Total & 7.55M & $\sim$9 hours & --- \\
\bottomrule
\end{tabular}
\end{table}

Inference cost for a single forecast at $N=100$ samples and $h=30$~min horizon: $\sim$0.3 seconds on GPU, $\sim$5 seconds on CPU. The score decoder is the bottleneck (100 reverse diffusion steps per sample), but could be accelerated with DDIM sampling.

% ============================================================
\section{Reproducibility}
\label{app:reproducibility}
% ============================================================

All experiments were run with three random seeds (42, 123, 456). Code, configuration files, and trained model checkpoints are available at \url{[redacted for review]}. The CloudCV dataset is publicly available from the NREL Data Catalog (submission 248). BMS meteorological data is available from the NREL SRRL MIDC portal.

\end{document}
